\subsection{Construction of the Chern Weil Homomorphism}

Let us recall the definition of the connection and curvature form. Let $E \to M$ be a vector bundle with connection $\nabla$. For a local frame $e_1,\ldots,e_k$ of $E$ over $U \subseteq M$ define the \textbf{connection form} as the unique matrix valued $1$-form $$\Theta = (\Theta_{ij})_{i,j=1}^k \in \Omega^1(U,\K^{k\times k})$$ with
\[
\nabla e_j = \sum_{i=1}^{k} \Theta_{ij} \otimes e_i.
\]
Using the connection form we can compute the covariant derivative of any section $s \in \Omega^0(M,E)$ as follows. In the chosen frame write $s = s^je_j$ for smooth functions $s^j \in C^\infty(M,\K)$, then 
\[
\nabla s = ds^j \otimes e_j + \sum_{i=1}^{k} s^j \Theta_{ij} \otimes e_i = (d + \Theta)s,
\]
where $d$ acts by differentiation with respect to the local frame as written above and $\Theta$ acts by matrix multiplication in the coordinates of this frame.

The \textbf{curvature form} $\Omega \in \Omega^2(U,\K^{k\times k})$ is then defined as 
\[
\Omega := d \Theta + \Theta \wedge \Theta.
\]
A computation shows that $\Omega$ transforms nicely under a change of frame and thus yields a globally defined object $\Omega \in \Omega^2(M,\mathrm{End}(E))$. We sometimes write $\Omega = \Omega_{\nabla}$ to denote the dependence on the connection $\nabla$. The curvature form behaves well under pullback:

\begin{lemma}
    \label{chern_weil:lem:naturality_of_curvature_form}
    Let $E \to M$ be a vector bundle with connection $\nabla$ and $f \colon N \to M$ a smooth map. Then 
    \[
    f^\ast \Omega_{\nabla} = \Omega_{f^\ast \nabla},
    \]
    where $f^\ast\nabla$ denotes the pullback connection on $f^\ast E$.
\end{lemma}
\begin{proof}
    Let $e_1\ldots,e_k$ be a local frame, let $\Theta$ be the associated connection form and let $\tilde{e}_j := f^\ast e_j$ be the pullback frame. Then by definition of the pullback connection (cf. Proposition \ref{riem_mnf:prop:pullback_connections}) we compute for any tangent vector $v \in T_qN$
    \[
    (f^\ast\nabla)_v \tilde{e}_j = \nabla_{f_\ast v} e_j = \sum_{i=1}^{k} (\Theta_{ij})_{f(q)}(f_\ast v)\, e_i(f(q)) = \sum_{i=1}^{k}(f^\ast \Theta_{ij})_q(v) \,\tilde{e}_j(q).
    \]
    Thus the connection form on $f^\ast E$ with respect to the frame $\tilde{e}_1,\ldots,\tilde{e}_k$ is given by $f^\ast \Theta$. Since pullback commutes with the exterior derivative and the wedge product the statement now follows from the definition of the curvature form.
\end{proof}

Chern-Weil theory now constructs cohomology classes (so called characteristic classes) of $M$ out of the curvature form using certain invariant homogeneous polynomials, which we shall introduce next.

\begin{definition}
    Let $V$ be an $n$-dimensional (real or complex) vector space. An $r$-homogeneous polynomial on $V$ is a map $\eta \colon V \to \K$, such that for all vectors $e_1,\ldots,e_n \in V$ the map 
    \[
    \K^n \to \K, \; (t_1,\ldots,t_n) \mapsto \eta(t_1e_1 + \ldots + t_n e_n)
    \]
    is a (multivariate) $r$-homogeneous polynomial in the variables $t_1,\ldots,t_n$. The space of all $r$-homogeneous polynomials on $V$ is denoted by $\mathrm{Pol}_r(V)$.
\end{definition}

By choosing $e_1,\ldots,e_n$ as a basis, it is easy to see that a map $\eta \colon V \to \K$ is an $r$-homogeneous polynomial if and only if for one (then all) linear isomorphisms $A \colon \K^n \to V$, the compostition $\eta \circ A \colon \K^n \to \K$ is an $r$-homogeneous polynomial. Thus we can think of such polynomials as $r$-homogeneous polynomials in the entries of a vector $v = (v_1,\ldots,v_n)$ with respect to some (then all) basis. With this definition it is easy to construct examples:

\begin{example}
    The determinant $\mathrm{det} \colon \K^{n\times n} \to \K$ is an $n$-homogenous polynomial and the trace $\mathrm{tr} \colon \K^{n\times n} \to \K$ is a $1$-homogeneous polynomial.
\end{example}

Yet another useful characterization of $\mathrm{Pol}_r(V)$, which we will make use of later, can be given as follows. 

\begin{lemma}
    Let $\Sigma^r_\K(V^\ast)$ be the space of $r$-linear symmetric maps $V \times \ldots \times V \to \K$. Then 
\begin{align*}
\Phi_r \colon \Sigma^r_\K(V^\ast) &\to \mathrm{Pol}_r(V) \\
    \overline{\eta} &\mapsto (v \mapsto \overline{\eta}(v,\ldots,v))
\end{align*}
is an isomorphism of vector spaces.
\end{lemma} 
\begin{proof}
    This can be seen using a polarisation argument: If $\eta = \Phi_r(\overline{\eta})$ then 
\[
\eta(t_1v_1 + \ldots + t_rv_r) = \overline{\eta}\left(\sum_{i=1}^{r} t_iv_i,\ldots,\sum_{i=1}^{r} t_iv_i\right)
\]
for vectors $v_i \in V$ and scalars $t_i \in \K$. Using multilinearity and symmetry of $\overline{\eta}$ we see that this is a polynomial in $t_1,\ldots,t_r$, where the coefficient in front of $t_1\cdots t_r$ is given by $r! \, \overline{\eta}(v_1,\ldots,v_r)$. So if $\eta = 0$ we conclude that $\overline{\eta}(v_1,\ldots,v_r) = 0$. Hence $\Phi_r$ is injective. To show surjectivity we argue similiarly. Let $\eta \in \mathrm{Pol}_r(V)$ and let $e_1,\ldots,e_n \in V$ be a basis. Then $\eta(t_1e_1 + \ldots + t_ne_n)$ is a polynomial in $t_1,\ldots,t_n$. Let $a_{i_1,\ldots,i_r}$ be the coefficient of this polynomial in front of $t_{i_1}\cdots t_{i_r}$, where we assume that $i_1 \leq \ldots \leq i_r.$ As above we see by comparing coefficients, that if $\eta = \Phi_r(\overline{\eta})$ then
\[
\overline{\eta}(e_{i_1},\ldots,e_{i_r}) = \frac{1}{r!} a_{i_1,\ldots,i_r}.
\]  
So we are led to define $\overline{\eta}$ on ordered tuples of basis vectors by this formula and extend it uniquely to a symmetric $r$-linear map. Then by definition the coefficients of the polynomials $\eta(t_1e_1 + \ldots + t_ne_n)$ and 
\[
\overline{\eta}\left(\sum_{i=1}^{n} t_ie_i,\ldots,\sum_{i=1}^{n} t_ie_i\right)
\]
agree, thus they are equal. Since $e_1,\ldots,e_n$ is a basis, this shows 
\[
\eta(v) = \overline{\eta}(v,\ldots,v).
\]
\end{proof}

We now turn our attention to homogeneous polynomials on the vector space $\K^{k\times k}$ of matrices, which are used to define characteristic classes in Chern-Weil theory. 

\begin{definition}
    An $r$-homogeneous polynomial $\eta \colon \K^{k \times k} \to \K$ is called invariant if $$\eta(A^{-1}BA) = \eta (B)$$ for all $B \in \K^{k\times k}$ and $A \in \mathrm{GL}(k,\K)$. We denote the space of all such polynomials by $I_r(k) \subseteq \mathrm{Pol}_r(\K^{k \times k})$.
\end{definition}

Under the isomorphism $\Phi_r$ the invariant polynomials correspond to the invariant $r$-linear symmetric maps, where an $r$-linear symmetric map $\overline{\eta} \colon \K^{k \times k} \times \ldots \times \K^{k \times k} \to \K$ is called invariant if 
\[
\overline{\eta}(A^{-1}B_1A,\ldots,A^{-1}B_rA) = \overline{\eta}(B_1,\ldots,B_r)
\]
for all $B_1,\ldots,B_r \in \K^{k \times k}$ and $A \in \mathrm{GL}(k,\K)$. Denote the space of all such maps by $S_r(k)$. We can then extend $\overline{\eta}$ to an $r$-linear map (here we don't need invariance yet)
\[
\overline{\eta} \colon \Omega^{l_1}(M,\K^{k \times k}) \times \ldots \times \Omega^{l_r}(M,\K^{k \times k}) \to \Omega^{l}(M,\K),
\]
where $l = l_1 + \ldots + l_r$ by setting 
\[
\overline{\eta}(\Omega_1,\scalebox{0.7}[1]{\ldots},\Omega_r)(X_1,\scalebox{0.7}[1]{\ldots},X_l) := \frac{1}{l_1! \scalebox{0.7}[1]{$\cdots$} l_r!} \sum_{\sigma} \mathrm{sgn}(\sigma) \overline{\eta}(\Omega_1(X_{\sigma(1)}, \scalebox{0.7}[1]{\ldots} ,X_{\sigma(l_1)}),\scalebox{0.7}[1]{\ldots}, \Omega_r(X_{\sigma(l-l_r + 1)}, \scalebox{0.7}[1]{\ldots} ,X_{\sigma(l)}))
\]
where the sum is taken over all permutations $\sigma$ of $l$ elements. The coefficient in the above definition is chosen, such that if $\Omega_i = \omega_i \otimes A_i$ for some $\omega_i \in \Omega^{l_i}(M,\K)$ and $A_i \in \Omega^0(M,\K^{k \times k})$ then
\[
\overline{\eta}(\Omega_1,\ldots,\Omega_r) = \omega_1 \wedge \ldots \wedge \omega_r \, \overline{\eta}(A_1,\ldots,A_r).
\]
So this definition agrees with the one in \cite{emilia_globale_ana}, but has the benefit of being coordinate independent. We will mostly work with this formula. Note that $\Omega_i$ can always be written as a sum of differential forms of the required form, where the $A_i$ are matrices with only one non-zero entry being equal to $1$. If $\Omega_1 = \ldots = \Omega_r =: \Omega$ and $\Omega$ has even degree, one can also think of $\overline{\eta}(\Omega,\ldots,\Omega)$ as the $l$-form obtained by evaluating the polynomial $\eta = \Phi_r(\overline{\eta})$ on the differential form valued entries of the matrix $\Omega$, where multiplication is given by the wedge product (the even degree of $\Omega$ is needed for the wedge product to be commutative).

Now given a vector bundle $E \to M$, the invariance of $\overline{\eta}$  allows us to define an $r$-linear map 
\[
\overline{\eta} \colon \Omega^{l_1}(M,\mathrm{End}(E)) \times \ldots \times \Omega^{l_r}(M,\mathrm{End}(E)) \to \Omega^{l}(M,\K)
\]
as follows: Choose any local frame of $E$ over $U \subseteq M$. Then under this frame $\Omega_i \in \Omega^{l_i}(M,\mathrm{End}(E))$ corresponds to a matrix valued form $\tilde{\Omega}_i \in \Omega^{l_i}(U,\R^{k \times k})$ and we set 
\[
\overline{\eta}(\Omega_1,\ldots,\Omega_r) := \overline{\eta}(\tilde{\Omega}_1,\ldots,\tilde{\Omega}_r)
\]
on $U$. Since $\overline{\eta}$ is invariant, this definition does not depend on the chosen frame, hence we get a well-defined global $l$-form. For a polynomial $\eta = \Phi_r(\overline{\eta})$ we write 
\[
\eta(\Omega) := \overline{\eta}(\Omega,\ldots,\Omega).
\]
We can now state the central theorem.

\begin{theorem}[Chern-Weil]
    \label{chern_weil:thm:chern_weil_thm}
    Let $E \to M$ be a rank $k$ vector bundle with connection $\nabla$ and curvature form $\Omega$ and let $\eta \in I_r(k)$. Then the following holds 
    \begin{enumerate}[(i)]
        \item $\eta(\Omega)$ is closed 
        \item If $\nabla^\prime$ is another connection with curvature form $\Omega^\prime$, then $[\eta(\Omega)] = [\eta(\Omega)]$
        in $H_{dR}^{2r}(M,\K)$.
    \end{enumerate}
\end{theorem}

We can thus define a map 
\[
\chi_E^r \colon I_r(k) \to H_{dR}^{2r}(M,\K)
\]
by choosing any connection on $E$ and defining $\chi_E^r(\eta) := [\eta(\Omega)]$ for the corresponding curvature form $\Omega$. In fact if we define the algebra 
\[
I(k) := \bigoplus_{r=0}^\infty I_r(k),
\]
where multiplication is defined pointwise as $(\eta_0 \cdot \eta_1)(A):= \eta_0(A)\eta_1(A)$, the maps $\chi_E^r$ define a homomorphism of commutative algebras
\[  
\chi_E \colon I(k) \to H_{dR}^{\text{even}}(M,\K),
\]
called the \textbf{Chern-Weil homomorphism}. This homomorphism satisfies the following naturality property.

\begin{corollary}
    Let $E \to M$ be a vector bundle and $f \colon N \to M$ a smooth map. Then 
    \[
    \chi_{f^\ast E} = f^\ast \chi_E.
    \]
\end{corollary}
\begin{proof}
    From the definition of $\eta(\Omega)$ it follows directly, that $\eta(f^\ast \Omega) = f^\ast\eta(\Omega)$, such that the statement follows together with Lemma \ref{chern_weil:lem:naturality_of_curvature_form}.
\end{proof}